\section{Discussion and Conclusion}
\label{sec:conclusion}

Training loss can fall in both constructive and destructive gradient regimes. Around the Adam loss spikes we traced, successive updates repeatedly moved predictions back and forth along a few functional directions. Capping the direct changes that Adam proposed at selected hidden sites reduced mean training loss and accumulated upward loss variation in all three models. The computational cost was substantial in the MLP, where evaluating the cap took longer than computing the full gradient. A practical strategy to test is to apply the cap only when the inexpensive gradient condition in Eq.~\eqref{eq:loss-rise-alignment}, $\rho_t<-\|g_{t-1}\|/\|g_t\|$, warns of a possible loss increase on the next update.

We chose modest-sized models and full-batch trajectories to connect the observed dynamics to the theory. We varied batch size to study gradient alignment, but traced representations and tested interventions under full-batch Adam with one initialization per model. Whether the same diagnostics and interventions remain effective in larger models under stochastic training is open. The cap reduced the largest output changes, yet functional oscillation persisted. Under q25, the Transformer's predictions still reversed direction on nearly nine out of ten updates.

The displacement ratios reveal a geometric inefficiency in both full-batch GD and Adam. Much of the distance traveled by GD in parameter space cancels across updates. In the Adam spike windows, net output displacement was only 1.7--3.5\% of the distance traveled, and most of that movement followed a few functional directions despite the number of parameters available. We propose testing a hybrid objective that combines the task loss with two auxiliary terms. One would encourage sparse or compact representation changes to limit oscillation along dominant directions. The other would increase the sensitivity of persistently inactive units to their inputs. The question is whether these objectives can reduce cancellation while bringing previously inactive units into learning.
